\documentclass[12pt]{article}
\usepackage[top=0.75in, bottom=1in, left=1in, right=1in]{geometry}
\usepackage{hyperref}
\usepackage[table]{xcolor}
\hypersetup{colorlinks=true, urlcolor=blue, linkcolor=blue}

\begin{document}

\vspace*{-0.8cm}
\noindent Prof. Eren Bilen
\hfill
\textbf{DATA 101 – Study Guide for Exam II}

\bigskip
\hrule
\bigskip

\noindent
\textbf{Overview:} \\[2pt]
The second exam focuses on geospatial analysis, modeling, supervised vs unsupervised learning, linear regression, model fit, and residuals.
\textit{This guide summarizes major ideas but does not list every detail covered in class.}


\bigskip
\hrule
\bigskip

% ------------------------------------------------------------
\noindent
\textbf{Maps in Python} \\[4pt]
Understand what shapefiles are and how they represent geospatial data.
Understand how to load geospatial data using GeoPandas.\\
Recognize how to filter data to a specific country or region.\\
Interpret simple static maps created with Matplotlib.\\
Understand how to add points (e.g., cities) using latitude and longitude.\\
Recognize how labels are added to maps using loops.\\
Understand the difference between static maps and interactive maps.\\
Recognize basic components of interactive maps created with Folium (map center, markers, tooltips).\\
\textbf{Keywords:} shapefiles, GeoPandas, Matplotlib, Folium, latitude, longitude, map, markers

\bigskip
\hrule
\bigskip


% ------------------------------------------------------------
\noindent
\textbf{Supervised vs Unsupervised Learning} \\[4pt]
Difference between predicting an outcome (supervised) and finding patterns or groups (unsupervised).\\
Examples of each: regression, classification, clustering.\\
Identify whether a task is classification or regression.\\
Recognize what makes a problem supervised: the presence of a response variable \(y\).\\
\textbf{Keywords:} supervised, unsupervised, prediction, clustering, classification, regression

\bigskip
\hrule
\bigskip


% ------------------------------------------------------------
\noindent
\textbf{Linear and Multilinear Regression} \\[4pt]
Understand the functional form \(y = a + bx\) and \(y = b_0 + b_1 x_1 + b_2 x_2 + \dots\).\\
Interpret slopes and intercepts in plain language.\\
Explain what a linear relationship means and when it might fail (curved patterns).\\
Recognize how multiple predictors affect interpretation.\\
Understand what \(R^2\) communicates about model fit.\\
\textbf{Keywords:} slope, intercept, predictor, response, linear relationship, \(R^2\)

\bigskip
\hrule
\bigskip


% ------------------------------------------------------------
\noindent
\textbf{Residuals and Model Fit} \\[4pt]
Define residuals as actual minus predicted values.\\
Interpret positive vs negative residuals.\\
Understand the goal of OLS: minimizing the sum of squared residuals.\\
Recognize good vs bad residual patterns (random cloud vs curvature).\\
Explain underprediction and overprediction using residuals.\\
\textbf{Keywords:} residual, prediction error, OLS, squared residuals, model fit

\bigskip
\hrule
\bigskip

\clearpage

% ------------------------------------------------------------
\noindent
\textbf{Training vs Testing} \\[4pt]
Understand why we split data: generalization to unseen cases.\\
Interpret differences between training error and testing error.\\
%Recognize overfitting and underfitting conceptually.\\
Explain why training data cannot be used for testing.\\
\textbf{Keywords:} train, test, generalization, overfitting, underfitting

\bigskip
\hrule
\bigskip

% ------------------------------------------------------------
\noindent
\textbf{Coding with Codex} \\[4pt]
Be able to read and interpret simple Python code generated by Codex.\\
Identify and fix basic coding errors (e.g., incorrect function names, syntax issues).\\
Understand that Codex can make mistakes and outputs must be checked carefully.\\
Interpret simple visualizations (bar charts with labels and titles).\\
\textbf{Keywords:} Codex, debugging

\bigskip
\hrule
\bigskip


\end{document}
